% Part: model-theory
% Chapter: models-of-arithmetic
% Section: introduction

\documentclass[../../../include/open-logic-section]{subfiles}

\begin{document}

\olfileid{mod}{mar}{int}
\section{परिचय}

अंकगणिताचे \emph{मानक प्रतिमान} म्हणजे अशी !!{structure}~$\Struct{N}$ की
$\Domain{N} = \Nat$ आणि तिच्यात $\Obj{0}$, $\prime$, $+$, $\times$ व
$<$ यांचे अपेक्षित नेहमीच्या अर्थाने अर्थनिर्धारण केलेले असते. म्हणजे,
$\Obj{0}$ हे $0$ असते, $\prime$ हे उत्तरवर्ती फलन असते, $+$ चा अर्थ
बेरीज आणि $\times$ चा अर्थ~$\Nat$ मधील संख्यांचा गुणाकार असा असतो.
अधिक स्पष्टपणे,
\begin{align*}
  \Assign{\Obj{0}}{N} & = 0\\
  \Assign{\prime}{N}(n) & = n + 1\\
  \Assign{+}{N}(n, m) & = n + m\\
  \Assign{\times}{N}(n, m) & = nm
\end{align*}
अर्थात, $\Lang{L_A}$ साठी अशाही रचना असतात की त्यांचे प्रांत~$\Nat$
नसतात. उदाहरणार्थ, $\Domain{M} = \{a\}^*$ (एकाच $a$ या चिन्हाच्या
सांत क्रमिका, म्हणजे $\emptyset$, $a$, $aa$, $aaa$, \dots) असा प्रांत
असलेली $\Struct{M}$ घेता येते आणि पुढील अर्थनिर्धारण करता येते:
\begin{align*}
  \Assign{\Obj{0}}{M} & = \emptyset\\
  \Assign{\prime}{M}(s) & = s \concat a\\
  \Assign{+}{M}(n, m) & = a^{n + m}\\
  \Assign{\times}{M}(n, m) & = a^{nm}
\end{align*}
या दोन रचना ``मूलतः एकच'' आहेत: फरक फक्त !!{domain}s चे
!!{element}s कोणते आहेत यात आहे; अर्थनिर्धारण फलनांमुळे !!{domain}s चे
!!{element}s परस्पर कसे संबंधित आहेत यात नाही. अशा दोन !!{structure}s ना
\emph{समरूप} म्हणतात.

संहतता प्रमेयावरून सहज निष्पन्न होते की~$\Struct{N}$ मध्ये सत्य असलेल्या
कोणत्याही उपपत्तीला~$\Struct{N}$ शी समरूप नसलेली प्रतिमानेही असतात. अशा
रचनांना \emph{अमानक} म्हणतात. त्यांच्याबद्दलची रंजक गोष्ट अशी की मानक
प्रतिमानातील !!{element}s (म्हणजे $\Struct{N}$ मधील, तसेच तिच्याशी समरूप
असलेल्या सर्व !!{structure}s मधील घटक) मानक संख्यांक~$\num{n}$ यांच्या
मूल्यांनी पूर्णपणे व्यापलेले असतात, म्हणजे
\[
\Domain{N} = \Setabs{\Value{\num{n}}{N}}{n \in \Nat}
\]
पण अमानक प्रतिमानांत तसे नसते: $\Struct{M}$ अमानक असेल, तर किमान एक
$x \in \Domain{M}$ असा असतो की प्रत्येक~$n$ साठी
$x \neq \Value{\num{n}}{M}$.

हे अमानक घटक फार रंजक आहेत: ते ``अनंत नैसर्गिक संख्या'' आहेत. पण त्यांच्या
अस्तित्वामुळे अपूर्णता-घटनांचे एका अर्थी स्पष्टीकरणही मिळते. उदाहरणार्थ,
पेआनो अंकगणितासाठीचे सुसंगतता-विधान $\OCon[\Th{PA}]$, म्हणजे $\lnot
\lexists[x][\OPrf[\Th{PA}](x, \gn{\lfalse})]$, विचारात घ्या.
$\Th{PA}$ कडून $\OCon[\Th{PA}]$ किंवा $\lnot \OCon[\Th{PA}]$ यांपैकी
काहीही सिद्ध होत नसल्यामुळे यांपैकी कोणतेही विधान~$\Th{PA}$ मध्ये सुसंगतपणे
जोडता येते. $\Th{PA}$ सुसंगत असल्यामुळे
$\Sat{N}{\OCon[\Th{PA}]}$, आणि म्हणून $\Sat/{N}{\lnot
  \OCon[\Th{PA}]}$. त्यामुळे $\Struct{N}$ हे $\Th{PA}
\cup \{\lnot \OCon[\Th{PA}]\}$ चे \emph{प्रतिमान नाही}, आणि त्या उपपत्तीची
सर्व प्रतिमाने अमानकच असली पाहिजेत. $\Th{PA} \cup \{\lnot
\OCon[\Th{PA}]\}$ च्या प्रतिमानात असा एखादा !!{element} असलाच पाहिजे की
तो $\lexists[x][\OPrf[\Th{PA}](x, \gn{\lfalse})]$ सत्य करणारा साक्षी
ठरेल, म्हणजे~$\Th{PA}$ पासून व्याघाताच्या !!a{derivation} ची गोडेल संख्या.
%READERNOTE{OLFOL-029 — गोठवलेल्या इंग्रजीतील या अस्तित्ववाचक पुराव्याच्या विधेयात साक्षीदार x हा पहिला युक्तिवाद चुकून वगळला आहे. त्याच परिच्छेदातील सुसंगतता-विधान आणि पुढील प्रत्येक मानक संख्यांकाचे सूत्र या दोन्हींनुसार मराठी सूत्रात x पुनःस्थापित केला आहे.}
असा !!a{element} मानक असू शकत नाही---कारण प्रत्येक~$n$ साठी
$\Th{PA} \Proves \lnot \OPrf[\Th{PA}](\num{n}, \gn{\lfalse})$.

\end{document}
